% ch5_answers.tex -- Module 6, Chapter 5 Quick check answers
\begin{chapteranswers}{5}{Applications of the Motor Effect: Motors, Generators and Braking}
  % 5.1 Torque on a coil and the simple DC motor
  \item Side: $\qty{2.5}{cm} = \qty{0.025}{m}$. Area:
        $A = 0.025 \times 0.025 = \qty{6.25e-4}{m^2}$.
        (a) Maximum torque: plane parallel to $B$, so $\theta = \ang{90}$.
        $\tau = nIAB = 200 \times 0.50 \times \num{6.25e-4} \times 0.30 = \qty{0.01875}{N.m}$.
        (b) The plane makes \ang{60} with the field, so
        $\theta = \ang{90} - \ang{60} = \ang{30}$.
        $\tau = nIAB\sin\theta = 0.01875 \times \sin\ang{30} = 0.01875 \times 0.50
        = \qty{0.009375}{N.m}$.
        Answers: \final{(a) $\qty{1.9e-2}{N.m}$} \ \final{(b) $\qty{9.4e-3}{N.m}$}.
  \item (a) Reversing the supply reverses the current in both long sides. By the
        right-hand palm rule, both forces reverse: the left side is now pushed up and
        the right side down. Answer: \final{clockwise} (seen from the commutator end).
        (b) Swapping the magnets reverses $B$, so both forces reverse again. Answer:
        \final{anticlockwise}, as in the worked example. Reversing \emph{both} the
        current and the field leaves the direction of turning the same.
  \item (a) Only one side is in the field, so find the force on that side first. The
        wire is at \ang{90} to $B$, so
        $F = lIB\sin\ang{90} = 0.040 \times 3.0 \times 0.050 = \qty{6.0e-3}{N}$.
        Then $\tau = r_\perp F = 0.15 \times \num{6.0e-3}$.
        Answer: \final{$\tau = \qty{9.0e-4}{N.m}$}.
        (b) $\tau = r_\perp lIB$, so $\tau \propto IB$. We need
        $\dfrac{\tau_2}{\tau_1} = \dfrac{I_2}{I_1} \times \dfrac{B_2}{B_1}$, so
        $2 = 4 \times \dfrac{B_2}{B_1}$ and $\dfrac{B_2}{B_1} = \dfrac{2}{4} = 0.5$.
        Answer: \final{$B$ must be halved}.
  \item (a) split-ring commutator; (b) brushes; (c) stator (the field magnets).
  \item The torque turns the coil towards the position where its plane is
        perpendicular to the field. The coil is moving, so it swings past this
        position. The torque then reverses and pulls it back. The coil rocks back and
        forth about that position, and friction makes it stop there. It never makes
        a full turn, because nothing reverses the current.
  % 5.2 Back emf and energy
  \item (a) At switch-on $\varepsilon_\mathrm{back} = 0$, so
        $I = \dfrac{V}{R} = \dfrac{9.0}{3.0} = \qty{3.0}{A}$.
        (b) $\varepsilon_\mathrm{back} = V - IR = 9.0 - 0.40 \times 3.0 = 9.0 - 1.2$.
        Answers: \final{$I = \qty{3.0}{A}$} and
        \final{$\varepsilon_\mathrm{back} = \qty{7.8}{V}$}.
  \item Vertical axis: current $I$ (\si{A}). Horizontal axis: time from switch-on
        (\si{s}). The graph starts at its largest value, $\dfrac{V}{R}$, at $t = 0$,
        because there is no back emf yet. It falls quickly at first, then more slowly.
        So the graph is a curve, not a straight line. This happens because the back emf
        grows as the motor speeds up. It levels off at a small value \emph{above zero},
        because friction (and any load) still needs some torque. (For an ideal motor it
        would fall towards zero.)
  \item (a) $\varepsilon_\mathrm{back} = V - IR = 18.0 - 3.0 \times 0.40 = \qty{16.8}{V}$.
        (b) Jammed: no rotation, so $\varepsilon_\mathrm{back} = 0$ and
        $I = \dfrac{18.0}{0.40} = \qty{45}{A}$.
        (c) Normal use: $P_\mathrm{heat} = I^2R = 3.0^2 \times 0.40 = \qty{3.6}{W}$.
        Jammed: $P_\mathrm{heat} = 45^2 \times 0.40 = \qty{8.1e2}{W}$, which is 225 times
        more. With no back emf, all of the input energy becomes heat in the coil. The
        wires overheat, the insulation melts and the motor can burn out.
        Answers: \final{\qty{16.8}{V}}, \final{\qty{45}{A}},
        \final{\qty{3.6}{W} and \qty{8.1e2}{W}}.
  \item Suppose the induced emf helped the supply instead. As the coil sped up, the
        emf would grow, so the current and the torque would grow. The coil would spin
        faster and faster with no limit. Its kinetic energy and the heating would grow
        without limit, more than the supply provides. Energy would be made from nothing,
        which breaks the law of conservation of energy. So the induced (back) emf must
        oppose the supply. This is Lenz's law.
  % 5.3 Generators
  \item An AC generator has two slip rings (full rings). Each end of the coil always
        stays joined to the same brush, so the output reverses every half turn (AC).
        A DC generator has a split-ring commutator. It swaps the connections every half
        turn, when the emf is zero, so the output always flows one way (pulsing DC).
  \item The emf depends on how fast the flux \emph{changes} (Faraday's law), not on the
        size of the flux. When the plane of the coil is parallel to the field, the
        flux is zero, but the sides cut straight across the field lines. The flux is
        changing fastest, so the emf is largest. (When the plane is perpendicular,
        the flux is largest but not changing for an instant, so the emf is zero.)
  \item By Lenz's law, the force on the induced current opposes the motion. The side
        moves down, so the force on it must point up. Palm rule: fingers along $B$ (to
        the right) and the palm pushing up. Then the thumb points out of the page. So
        the induced current in the right-hand side flows \final{out of the page}
        (towards you).
  \item Start (plane perpendicular to $B$, $\theta = 0$):
        $\Phi_\mathrm{initial} = BA = 0.30 \times \num{1.2e-3} = \qty{3.6e-4}{Wb}$.
        End (plane parallel, $\theta = \ang{90}$): $\Phi_\mathrm{final} = 0$.
        $\Delta\Phi = 0 - \num{3.6e-4} = \qty{-3.6e-4}{Wb}$.
        $\varepsilon = -N\dfrac{\Delta\Phi}{\Delta t}
        = -50 \times \dfrac{\num{-3.6e-4}}{0.020} = \qty{0.90}{V}$.
        Answer: \final{average emf $= \qty{0.90}{V}$}.
  \item (a) $f = \dfrac{1500}{60} = \qty{25}{Hz}$, so
        $T = \dfrac{1}{f} = \dfrac{1}{25} = \qty{0.040}{s}$.
        (b) Half the speed: the flux changes half as fast, so the peak halves:
        $8.0 \times \dfrac{750}{1500} = \qty{4.0}{V}$. Each turn takes twice as long:
        $f = \dfrac{750}{60} = \qty{12.5}{Hz}$ and $T = \dfrac{1}{12.5} = \qty{0.080}{s}$.
        Answers: \final{(a) \qty{0.040}{s}} \ \final{(b) \qty{4.0}{V}, \qty{0.080}{s}}.
  % 5.4 AC induction motors and magnetic braking
  \item The rotor currents are induced by the changing flux of the rotating field
        (Faraday's law). No current has to be fed into the rotor, so no sliding
        contacts (brushes) or commutator are needed. The stator coils do not move,
        so they can be wired straight to the supply.
  \item (a) Slip $= 3000 - 2910 = \qty{90}{rpm}$. In turns each second:
        $90 \div 60 = 1.5$. Answer: \final{\qty{90}{rpm} $=$ 1.5 turns each second}.
        (b) A bigger load slows the rotor, so the slip increases. The field sweeps past
        the bars faster, so the flux through the bar loops changes faster. The induced
        emf and current are larger, so the forces on the bars and the torque are larger.
        The motor settles at a slightly lower speed, where its torque balances the load.
  \item The braking force comes from eddy currents. These need a changing flux, and the
        flux changes only while the metal moves through the field. When the train
        stops, there is no flux change, no eddy current and no force. The force also
        becomes weak at low speed. So friction brakes are needed to stop and hold it.
  \item $Q = \tfrac{1}{2}mu^2 - \tfrac{1}{2}mv^2
        = \tfrac{1}{2}(\num{3.0e4})(20^2) - \tfrac{1}{2}(\num{3.0e4})(12^2)
        = \num{6.0e6} - \num{2.16e6} = \qty{3.84e6}{J}$.
        Answer: \final{$Q = \qty{3.8e6}{J}$}.
  \item The disc moves through the magnet's field, so the flux through parts of the
        disc changes. By Faraday's law, emfs are induced and eddy currents flow. By
        Lenz's law, the eddy currents make a force that opposes the motion, so the disc
        slows. Its kinetic energy becomes electrical energy and then heat in the disc.
        If the force helped the motion, the disc would speed up by itself and gain
        energy from nothing. That would break the law of conservation of energy, so
        Lenz's law must be true.
\end{chapteranswers}
